\documentclass[12pt,letterpaper,oneside]{article}
\input{ITESCA/maestria-en-gestion-administrativa/plan-de-negocios-mga/actividades-generadas/2026-II/revision-2026-09-20/formato-itesca-plan-negocios.tex}
\begin{document}
\portadaITESCA{Resultados de la investigación de mercados}{6549}
\textbf{Borrador hipotético pendiente de confirmación.} Este documento presenta una estructura de análisis para comunicar los resultados de la investigación de mercados de AM Taller Autocentro. Debido a que todavía no se dispone del trabajo de campo, la base de respuestas ni la identificación verificada de oferentes, no se reportan cifras, porcentajes, participantes, competidores o conclusiones sobre la demanda. El contenido deberá actualizarse con evidencia real antes de considerarse un resultado definitivo.

\section{Propósito de la investigación}

La investigación tiene como propósito reunir evidencia para valorar la demanda potencial de los servicios propuestos por AM Taller Autocentro, proyecto local concebido como taller, llantera y refaccionaria. El estudio deberá examinar necesidades relacionadas con mantenimiento y reparación mecánica, adquisición o instalación de llantas y búsqueda de refacciones. Asimismo, deberá identificar oferentes de productos o servicios similares y alternativas sustitutas dentro del ámbito geográfico que se confirme.

Los objetivos específicos son:

\begin{itemize}
    \item Describir las características relevantes de las personas consultadas y los criterios utilizados para seleccionarlas.
    \item Analizar necesidades, hábitos de búsqueda y factores de elección vinculados con la atención automotriz.
    \item Identificar indicios de demanda potencial mediante respuestas verificables.
    \item Reconocer oferentes directos, indirectos, especializados y sustitutos.
    \item Formular oportunidades y riesgos comerciales sin presentar las hipótesis del proyecto como hechos comprobados.
\end{itemize}

\section{Metodología y alcance pendiente}

La metodología deberá documentarse después de realizar el sondeo. Será necesario precisar el tipo de estudio, medio de aplicación, periodo de levantamiento, población de referencia, ámbito geográfico, criterios de inclusión y cantidad de respuestas válidas. También deberá explicarse si la selección de participantes fue probabilística o no probabilística, así como las limitaciones derivadas del procedimiento empleado.

Actualmente están pendientes el instrumento aplicado, el tamaño de la muestra, las fechas del trabajo de campo, el número de invitaciones, las respuestas recibidas, los registros descartados y los criterios de depuración. Si se utiliza una muestra por conveniencia, los resultados deberán presentarse como indicios exploratorios y no como estimaciones representativas de todo el mercado.

Antes del análisis será necesario revisar la integridad, pertinencia y consistencia de los registros. Las respuestas duplicadas, incompletas o ajenas al perfil definido deberán identificarse mediante criterios uniformes establecidos previamente. No deberán modificarse ni excluirse datos con el propósito de favorecer la viabilidad comercial del proyecto.

\section{Presentación de los resultados del sondeo}

\subsection{Descripción de participantes}

Esta sección deberá informar solamente las características incluidas en el instrumento y pertinentes para la investigación. Podrían considerarse el tipo de usuario de vehículo, la frecuencia aproximada de mantenimiento o la zona general de residencia, siempre que tales datos hayan sido solicitados de manera adecuada. Permanecen pendientes el número de participantes y su distribución real.

La versión final deberá distinguir entre el total de respuestas recibidas y el total de respuestas válidas. Cuando una pregunta admita omisiones o selecciones múltiples, será necesario explicar el total empleado como base de cálculo. La información personal deberá presentarse de manera agrupada para evitar la identificación de participantes.

\subsection{Necesidades y comportamiento de compra}

Los hallazgos podrán organizarse mediante los siguientes ejes, siempre que correspondan a preguntas efectivamente aplicadas:

\begin{enumerate}
    \item Necesidades de mantenimiento, reparación, llantas o refacciones señaladas por las personas participantes.
    \item Frecuencia con la que se busca atención automotriz o se adquieren insumos relacionados.
    \item Medios utilizados para localizar, consultar y comparar alternativas.
    \item Factores de elección, como confianza, claridad del diagnóstico, tiempo de atención, garantía, ubicación o disponibilidad.
    \item Dificultades experimentadas al solicitar servicios o buscar refacciones.
    \item Disposición declarada para conocer o solicitar una cotización de una nueva alternativa.
\end{enumerate}

Cada hallazgo deberá acompañarse de su frecuencia real y, cuando corresponda, de un porcentaje calculado sobre respuestas válidas. Las respuestas abiertas requerirán categorías transparentes y una explicación del procedimiento de clasificación. Los ejemplos textuales deberán anonimizarse y proceder de respuestas reales; no se crearán testimonios ni se atribuirán opiniones a participantes inexistentes.

\subsection{Interpretación de la demanda potencial}

La demanda potencial no puede sostenerse únicamente mediante expresiones generales de interés. Una interpretación favorable requeriría observar coincidencia entre necesidades concretas, frecuencia de consumo, dificultades con las alternativas actuales y disposición a solicitar información o una cotización. Incluso si aparecen estas señales, la conclusión deberá formularse como una estimación sujeta a validación.

El análisis distinguirá entre interés declarado, intención de compra y comportamiento efectivo. Si la cantidad y calidad de los datos lo permiten, podrán compararse segmentos con necesidades diferentes. No se asignarán precios, ventas esperadas ni volúmenes de demanda sin una base empírica específica.

\section{Análisis de oferentes y sustitutos}

La revisión competitiva deberá realizarse en el territorio confirmado para el proyecto. Por el momento no es posible nombrar empresas ni atribuirles precios, reputación, capacidades, inventarios o participación de mercado.

Los oferentes identificados podrán clasificarse de la siguiente manera:

\begin{itemize}
    \item \textbf{Directos:} establecimientos con una combinación comparable de taller, llantas o refacciones.
    \item \textbf{Especializados:} negocios enfocados en mecánica, diagnóstico, suspensión, llantas u otra categoría particular.
    \item \textbf{Indirectos:} comercios o prestadores que atiendan solamente una parte de la necesidad.
    \item \textbf{Sustitutos:} agencias, servicios móviles, compras digitales u otras alternativas verificadas que resuelvan la necesidad de una manera distinta.
\end{itemize}

Para cada oferente se deberá registrar el nombre comercial verificable, ubicación general, categoría, servicios anunciados, atributos observables, fuente y fecha de consulta. Los precios, existencias, promociones, compatibilidades y opiniones de clientes solo deberán incorporarse cuando exista evidencia documentada y comparable. La observación deberá aplicar criterios equivalentes a todos los oferentes para evitar comparaciones parciales.

\section{Conclusiones provisionales}

En ausencia de datos de campo, no es válido concluir que existe demanda suficiente o que AM Taller Autocentro posee una ventaja competitiva. La integración propuesta de taller, llantera y refaccionaria, junto con agenda de citas, control de inventario y revisión de proveedores, constituye una hipótesis que debe contrastarse con necesidades reales. No se afirma que estas funciones estén implementadas ni que el proyecto cuente con ventas, clientes, existencias o acuerdos comerciales.

La valoración de la viabilidad comercial dependerá de la consistencia de los resultados, el reconocimiento de sus limitaciones y la comparación verificable con oferentes y sustitutos. Si la evidencia resulta insuficiente o contradictoria, deberá ampliarse el sondeo, revisarse el segmento o ajustarse la propuesta antes de elaborar proyecciones comerciales.

\section{Datos necesarios para convertir el borrador en trabajo real}

Se requieren el instrumento definitivo, la base anonimizada de respuestas, los criterios de depuración, el tamaño y las características de la muestra, las fechas y el medio de levantamiento, el territorio analizado, los resultados por pregunta y la evidencia de los oferentes identificados. También permanecen pendientes el formato oficial, la rúbrica, el peso individual y las indicaciones complementarias de la actividad de GGPN01, Plan de Negocios, de la Maestría en Gestión Administrativa del ITESCA.
\clearpage
\printbibliography[title={Fuentes de consulta}]
\end{document}
